\section{Related Work}\label{sec:related}

\subsection{Relatedness, overlap and domain (RQ1)}\label{sec:rw-rq1}
Telmo Pires et al.~\cite{pires2019} showed that multilingual encoders transfer NER between languages with almost no lexical overlap; however, overlap still predicts how much transfer occurs. For African NER, Michael Beukman and Manuel Fokam~\cite{beukman2023} found that token overlap correlates with transfer at Pearson $R=0.73$, against $-0.30$ for genetic distance, but their ten languages include no Southern Bantu language. In MasakhaNER~2.0, geographic distance and entity overlap were the most important features of the LangRank source ranker~\cite{lin2019}, and entity overlap correlated with transfer at Spearman $R=0.6$~\cite{adelani2022}. Furthermore, embedding similarity predicts NER transfer but only weakly predicts sentiment transfer~\cite{idris2026}, which means that no single predictor holds across tasks.

Studies that separate language shift from domain shift often find that domain costs as much as language. In DaN+, Barbara Plank et al.~\cite{plank2020} report that German\ra Danish news lost about 4~F1, while same-language shifts from news to social media lost about 13--40, and Maarten De Raedt et al.~\cite{deraedt2023} found for sentiment that domain-only drops are at least as large as language drops. In contrast, Antonio Martínez-García et al.~\cite{martinezgarcia2021} use correlational domain proxies and conclude that most transfer loss is not due to domain. The domain shifts in CrossNER also add entity types~\cite{liu2021crossner}, which ours do not. For South African languages, MphayaNER is the only direct contrast we found: Rendani Mbuvha et al.~\cite{mbuvha2023} report that every related-language news source beat Tshivenda government data on Tshivenda news (best 46.2 against 35.2), although from one target and without significance tests. Government-domain zero-shot transfer also appears to be hard, since NGLUEni reports 32.4 weighted F1 for isiXhosa\ra isiZulu on SADiLaR NER data~\cite{meyer2024}, against 82.1 entity F1 for the same pair on news in MasakhaNER~2.0, with a different model and metric~\cite{adelani2022}. Thus, no study crosses source language with source corpus for NER among South African Bantu languages, and in NCHLT versus MasakhaNER the domain coincides with the annotation guideline, which none of these studies separates.

\subsection{Datasets with different annotation conventions}\label{sec:rw-guidelines}
Training one model on datasets annotated under different conventions has precedents. Sara Stymne et al.~\cite{stymne2018} improved a multi-treebank dependency parser with a learned treebank embedding, by up to 3.5 LAS over single-treebank training and by over 2.0 over concatenation. Similarly, Gamal Crichton et al.~\cite{crichton2017} found that multi-task biomedical NER with dataset-specific outputs helped small datasets the most, and taggers have been trained over datasets whose tag sets only partly overlap~\cite{greenberg2018,huang2019}. Bottleneck adapters~\cite{houlsby2019,pfeiffer2020} give each dataset its own small set of parameters on a shared encoder. However, we found no study that separates annotation conventions from text domain in cross-corpus NER for these languages, and Section~\ref{sec:future} proposes how to do so.

\subsection{Pretraining exposure (RQ2)}\label{sec:rw-rq2}
Wietse de Vries et al.~\cite{devries2022} fitted a mixed-effects model of zero-shot POS transfer over more than 100 languages and found that target inclusion in pretraining was the strongest predictor, with an estimated 19.2\% higher accuracy. However, the gains from adaptation point the other way: in Glot500, languages unseen by XLM-R gain $+13.2$~F1 on WikiANN from inclusion, against $+2.1$ for head languages~\cite{imani2023}, and adaptation benefits unseen languages more than seen ones~\cite{ansell2022,faisal2022}. Furthermore, Phillip Rust et al.~\cite{rust2021} found that tokenizer quality matters about as much as the size of the pretraining data, therefore an exposure effect may partly be a tokenizer effect. AfroXLMR's adaptation corpus contains isiZulu, isiXhosa and Sesotho but not siSwati, Sepedi or Setswana~\cite{alabi2022}. In NGLUEni, AfroXLMR-large improved on XLM-R-large by $+2.1$ for seen isiZulu but by $+5.3$ and $+5.9$ for unseen isiNdebele and siSwati~\cite{meyer2024}. Thus, no study reports seen-versus-unseen NER transfer ratios for these languages, and every unseen target has a seen close relative.

\subsection{Few-shot target data (RQ3)}\label{sec:rw-rq3}
Anne Lauscher et al.~\cite{lauscher2020} fine-tuned XLM-R on English and then on $k$ target sentences, and it averaged 85.43 F1 on WikiANN at $k=0$, 91.49 at $k=100$ and 93.69 at $k=500$, which means that most of the gain arrives by 100 sentences. Mengjie Zhao et al.~\cite{zhao2021} found that the choice of sampled sentences causes more variance than the seed, and that plain fine-tuning beats more elaborate few-shot methods. In MasakhaNER~2.0, the best zero-shot source comes close to annotating 500 target sentences in most cases~\cite{adelani2022}. However, no study compares the learning curves of a related-language source and a same-language source from another domain for NER.

\subsection{Sources for Setswana (RQ4)}\label{sec:rw-rq4}
In MasakhaNER~2.0, Setswana's best single sources were German (73.3) and Swahili (73.1), LangRank's predicted top two were Swahili and Chichewa~\cite{lin2019}, and the pool had no Sotho-Tswana language~\cite{adelani2022}. The most direct evidence that Sotho-Tswana languages help each other comes from speech, where Christiaan Jacobs and Herman Kamper~\cite{jacobs2021} found that acoustic word embeddings for Sesotho trained on Sepedi and Setswana reached 76.7\% AP, against 51.9\% for isiXhosa and siSwati. PuoBERTa evaluates Setswana NER without cross-lingual experiments~\cite{marivate2023}, and the transfer study of Sindane uses static embeddings~\cite{sindane2023}. Thus, we found no transformer-based NER study that uses Sepedi or Sesotho as a source for Setswana.
